\chapter{Of Assembly and the Common Halls}
\label{art:assembly}

\articlenote{In which is set down the Assembly of citizens, the common halls where the Assembly gathers, the speaking of citizens before the Assembly, and the will of the Assembly to amend this Pact or overturn the Mayor's or Judicial's acts where the Pact requires it.}

\section{The Assembly of Citizens}

The Assembly is the body of all citizens of majority, gathered in one place or by messenger-relay across the floors under Article 20, to hear the Mayor's report and the concerns of the silo, to confirm or remove the Mayor from office, to confirm or remove the Judge from office, to ratify amendments to this Pact, and to rule on disputes between the Mayor and Judicial where this Pact requires their judgment.

\begin{clauses}
\item The Mayor convenes the Assembly once yearly, at a time of the Mayor's choosing, to report on the silo's condition and to make recommendations to the citizens. Any citizen of majority may attend this ordinary Assembly and speak.
\item The Assembly may be convened in extraordinary session by petition of one fifth of the silo's population of majority, or by the Mayor, or by Judicial where this Pact requires the Assembly's judgment on a matter between the Mayor and Judicial. An extraordinary session must be called within fourteen days of the petition or circumstance requiring it.
\item No Assembly session shall last longer than three days without recess; a matter that requires deliberation longer than three days shall be put to the Assembly again, not in continuous session.
\item A citizen may send a message to the Assembly in their stead under Article 20, if the citizen is unable to attend in person; such a message may be read aloud in the Assembly and is recorded in the archive.
\end{clauses}

\section{The Common Halls and the Gathering}

The Assembly gathers in the large common halls on the Up Top, which are kept by the Department of Supply under Article 11. These halls are reserved for the Assembly and for such other gathering as the Mayor authorizes under Article 4.

\begin{clauses}
\item The main common hall of the Assembly, on Level 1, can hold several hundred citizens, and the Mayor shall ensure adequate seating and ventilation for the Assembly. Where more citizens wish to attend than the hall can hold, the overflow may gather at secondary halls and relay the proceedings by porter-relay or, where Information Technology under Article 8 has provided the means, by radio relay to other floors.
\item The Sheriff keeps order in the common halls where the Assembly gathers, and may remove a citizen who disrupts the Assembly, provided the citizen is given the reason for removal and may request a hearing before Judicial as to whether the removal was lawful.
\item A citizen making a spectacle or a threat in the Assembly halls is answerable under Article 17 as for public disruption or endangerment of the Assembly itself.
\item The common halls that serve as cafeterias open before the first labor of the day and close after the last, save the Down Deep's own hall, kept to the hours of Mechanical's and Mines's own shifts; a citizen whose labor keeps unusual hours may petition the relevant Head of Department for standing leave to draw food outside these hours, and such leave, once granted, need not be asked again.
\item Where the press of citizens wishing to witness a cleaning at the Level 1 wallscreen exceeds the hall's capacity, seating is granted by the same lottery Article 3 provides for the bearing of children, drawn separately and for this purpose alone, that the granting of a scarce place be seen to rest on chance and not on favor.
\end{clauses}

\section{The Hearing Before the Assembly}

Any citizen may petition to speak before the Assembly, and the Assembly shall hear petitions from any citizen who wishes to raise a matter of silo concern. The Mayor sets the time and order of speaking, but may not refuse a citizen the right to speak, save where the citizen has already disrupted an Assembly.

\begin{clauses}
\item A citizen wishing to speak must declare the subject of the address, and the Mayor may set limits on the time of speaking (no single citizen more than a few minutes; a matter of great complexity may be split across multiple speakers or multiple Assembly sessions).
\item A citizen's address to the Assembly is entered in the archive under Article 6 and may not be censored or removed, save where the address itself is sedition or forbidden speech under Article 16, in which case the Judge may review the archive copy and rule whether the address was rightfully censored.
\item A citizen who speaks to the Assembly is protected from retaliation under this Pact; any citizen punished by a Department, the Sheriff, or the Mayor for the fact of having spoken to the Assembly is answerable under Article 17 as for breach of the Pact.
\end{clauses}

\section{The Assembly's Powers of Confirmation and Removal}

The Assembly, together with the Mayor or Judicial as the case requires, confirms citizens to the highest offices and may remove them. The Mayor serves at the Assembly's pleasure, and the Judge serves at the Assembly's and Mayor's pleasure.

\begin{clauses}
\item The Assembly, in extraordinary session, votes on the removal of the Mayor by petition or by recommendation from Judicial; a removal vote passes upon the affirmative vote of two thirds of the Assembly present.
\item The Assembly, in extraordinary session and joined by the Mayor, votes on the confirmation of a new Judge or the removal of a sitting Judge; a confirmation requires a simple majority, and a removal requires two thirds.
\item The Assembly shall, when the Mayor convenes an ordinary session, receive a report from the Mayor on the major decisions, standing orders, and resource allocations of the past year, and may petition Judicial to rule on whether any such decision violated the Pact.
\end{clauses}

\section{The Amendment of This Pact}

This Pact may be amended only by the Assembly, upon petition from the Mayor or from one fifth of the citizens of majority, and the amendment must pass by two thirds of the Assembly present. Amendments are entered in the archive and in a separate roll of amendments under Article 23.

\begin{clauses}
\item The Mayor may recommend an amendment to this Pact and shall present the text of the proposed amendment to the Assembly for deliberation and vote.
\item Any citizen or group of citizens may petition for an amendment, provided the petition is signed by at least one fifth of the silo's population of majority, and the Assembly shall vote on the proposed amendment.
\item An amendment, once ratified by the Assembly, becomes part of this Pact and binds the silo from that date forward, exactly as if it had been written at the Pact's founding.
\item No amendment shall contradict a previous amendment or a prior Pact provision without explicitly repealing or replacing that provision; all amendments are entered in the roll and the archive, so the silo's citizens and Judicial may see how this Pact has grown and changed over time.
\end{clauses}
